\chapter{GamerServices and Local Identity}
\label{ch:gamerservices}

XNA tied \cnans{GamerServices} to Xbox Live. FNA retains much of the API surface on PC but
leaves the service-dependent paths inert. CNA supplies offline identities, achievements,
leaderboards, and two Guide overlays. It does not provide authentication, friends, matchmaking,
or a remote account service.

This chapter uses three statuses:

\begin{description}
  \item[Implemented] The API has local in-memory behavior.
  \item[Locally persisted] State survives through CNA's file store, with CNA-defined policy where
        Xbox Live supplied the original policy.
  \item[Stub] The API shape exists, but no local service can perform the operation.
\end{description}

These statuses describe the pinned revision; they are not claims of Xbox Live compatibility.

\section{Local identities and object lifetime}

\cnaclass{Gamer} is the base of \cnaclass{SignedInGamer}. Initialization creates four fixed
local identities: \texttt{Stub Gamer} and \texttt{Stub Gamer (1)} through \texttt{(3)}. Player
One is not a guest; the other three are guests. All four report
\texttt{IsSignedInToLive=true}, although no online sign-in occurs.

\texttt{GetFromGamertag}, \texttt{GetPartnerToken}, and their asynchronous pairs are stubs.
\texttt{IsFriend} returns \texttt{false}, and \texttt{GetFriends()} returns an empty
\cnaclass{FriendCollection}. The friend types themselves work as containers; CNA has no source
from which to populate them.

\begin{limitationnote}{Address stability}
A \cnaclass{Gamer} owns a \cnaclass{LeaderboardWriter} that stores the owning
\texttt{Gamer*}. The default copy and move operations copy that pointer. Moving a gamer after
construction can therefore leave its writer pointing at the old address. Keep gamers at stable
addresses---heap allocation is the documented choice---once their leaderboard writer may be
used.
\end{limitationnote}

XNA exposes \texttt{Presence} as a get-only reference property whose returned object is
mutable. CNA adds a non-\texttt{const} \texttt{getPresenceProperty()} overload so that the same
usage remains possible in C++:

\begin{lstlisting}
SignedInGamer* gamer = /* obtain from SignedInGamer::SignedIn */;
gamer->AwardAchievement("ACH_FIRST_BOSS_DEFEATED");
AchievementCollection earned = gamer->GetAchievements();

gamer->getPresenceProperty().setPresenceModeProperty(
    GamerPresenceMode::SinglePlayer);
\end{lstlisting}

The fragment assumes the corresponding GamerServices headers and namespace imports.

\section{Persistence contract}

The store is rooted at \cnaclass{StorageDevice}\texttt{::GetStorageRootEXT()}, which is based
on \texttt{SDL\_GetPrefPath}. Achievement files live below
\texttt{GamerServices/achievements}; leaderboard files live below
\texttt{GamerServices/leaderboards}. Calling
\texttt{StorageDevice::SetAppNameEXT} changes the root for both systems. Chapter~\ref{ch:storage}
defines that shared boundary.

Achievements use one JSON file per sanitized gamertag, with records of the form
\texttt{\{"key":...,"earnedTicks":...\}}. Leaderboards use one file per sanitized
\texttt{(key)\_(gameMode)}. Characters outside \texttt{[A-Za-z0-9.\_-]} become underscores;
an empty name becomes \texttt{\_}. Writes use an adjacent temporary file followed by rename,
with direct writing as a fallback. Missing or malformed files produce an empty store.

JSON numbers are represented as \texttt{double}. Achievement ticks, leaderboard ratings, and
64-bit property values therefore lose low bits above $2^{53}$. For a 2026
\cnaclass{DateTime}, the quantization step is about 128 ticks, or 12.8 microseconds. The stored
integer remains stable after the first serialization, but the original 64-bit value is not a
bit-exact round trip.

\section{Achievements}

\texttt{AwardAchievement} and \texttt{GetAchievements} use the local store. Returned entries
contain a persisted \texttt{Key}, \texttt{IsEarned=true}, and the earned time, subject to the
numeric limitation above. \texttt{Name}, \texttt{Description}, and \texttt{GamerScore} remain
empty or zero; \texttt{DisplayBeforeEarned} remains \texttt{true}. Xbox Live supplied that
catalog metadata outside the \texttt{AwardAchievement(string)} call, and CNA has no equivalent
catalog. Applications that display achievement names or scores must maintain their own mapping.

\texttt{Achievement::GetPicture()} throws \texttt{NotImplementedException}. Artwork is another
service-owned field for which the local store has no source.

\section{Leaderboards}
\label{sec:leaderboards}

\cnaclass{LeaderboardReader}, \cnaclass{LeaderboardWriter}, and
\cnaclass{LeaderboardEntry} are disk-backed. Their policy is CNA-specific because FNA exposes
no functioning PC implementation from which to inherit sorting or paging behavior:

\begin{itemize}
  \item entries sort by rating, descending;
  \item paging around an absent gamer starts at the top;
  \item persisted entries whose gamertags have no matching live \cnaclass{Gamer} are omitted;
  \item the asynchronous pairs complete synchronously;
  \item assigning \texttt{Rating} on a writer-created entry persists it immediately.
\end{itemize}

A reader-created entry has no persistence hook. Changing its in-memory rating does not update
the file.

\begin{lstlisting}
LeaderboardWriter& writer = gamer->getLeaderboardWriterProperty();
LeaderboardIdentity board = LeaderboardIdentity::Create(
    LeaderboardKey::BestScoreLifeTime);

LeaderboardEntry* entry = writer.GetLeaderboard(board);
entry->setRatingProperty(newHighScore); // the assignment is the commit
\end{lstlisting}

\section{Guide overlays and stubs}

Most \cnaclass{Guide} UI calls are no-ops: compose message, friend request, friends, game
invite, gamer card, marketplace, messages, party, player review, players, sign-in, and
\texttt{ShowAchievementsEXT}. \texttt{DelayNotifications} is also inert.

Two asynchronous families are CNA-rendered overlays:
\texttt{BeginShowMessageBox}/\texttt{EndShowMessageBox} and
\texttt{BeginShowKeyboardInput}/\texttt{EndShowKeyboardInput}. They stay pending until the
game's \texttt{Draw()} path invokes the matching \texttt{RenderPending*EXT} hook with a
\cnaclass{SpriteBatch}, \cnaclass{SpriteFont}, and white pixel. The hooks render the prompt and
poll keyboard or mouse input. \texttt{Guide.IsVisible} reports whether either operation is
pending.

Keyboard input preserves UTF-16 surrogate pairs. Password mode masks displayed code units but
returns the unmasked text. Escape cancels the operation. Because \texttt{std::string} cannot
distinguish XNA's null-on-cancel result from a confirmed empty string, callers must query
\texttt{WasKeyboardInputCanceledEXT} before interpreting the result.

\begin{lstlisting}
System::IAsyncResult* pending = Guide::BeginShowKeyboardInput(
    PlayerIndex::One, "Enter Passphrase", "Protect this save file",
    "", nullptr, std::any{}, true);

// Call once per Draw while the request is pending.
Guide::RenderPendingKeyboardInputEXT(
    *graphicsDevice, spriteBatch, uiFont, whitePixel);

if (!Guide::getIsVisibleProperty()) {
    if (!Guide::WasKeyboardInputCanceledEXT(pending)) {
        std::string passphrase = Guide::EndShowKeyboardInput(pending);
    }
    delete pending; // the caller owns the IAsyncResult
}
\end{lstlisting}

Only one Guide overlay can be pending. The test corpus covers pending state, callback identity,
Enter and Escape, prompt fields, password display, and competing requests.

\begin{historynote}{Why the overlay contract is stated explicitly}
An independent review found that an earlier implementation ignored the title and description,
reported \texttt{IsVisible=false}, displayed password text without masking, and had no cancel
path. The current implementation and tests cover all four corrections. The episode is useful
because a completed task record had overstated the observable behavior.
\end{historynote}

\section{Presence, privileges, and profile}

\cnaclass{GamerPresence} supplies all 60 XNA display strings and formats parameterized modes.
Its \texttt{SetPresenceModeStringEXT} publication hook is empty: no service consumes the
result. \cnaclass{GamerPrivileges} is fixed and permissive, and nothing enforces it.
\cnaclass{GamerProfile} uses synthetic defaults except for host-derived \texttt{Region}. These
objects carry useful local state, not authentication or policy authority.

\section{Dispatcher, component, and the historical hang}

\cnaclass{GamerServicesComponent} must be added to \texttt{Game.Components} before using the
Gamer or Guide APIs, as on XNA 4.0 for Windows. Its dispatcher creates and frees the local
gamer set. The \texttt{Update()} method on \cnaclass{GamerServicesDispatcher} is empty; the local asynchronous
implementations complete their own actions.

\label{sec:gamerservices-hang}
\begin{historynote}{Cross-namespace completion bug}
After initialization, \texttt{GamerServicesDispatcher::}\allowbreak\texttt{UpdateAsync()}
returns \texttt{true}
indefinitely. FNA's \cnaclass{NetworkSession} action logic waited for it to become
\texttt{false}; combining GamerServices initialization with network creation could therefore
spin forever. \texttt{SignedInGamer::BeginGetAchievements} once had the same dependency.

CNA's network create, find, and join work is performed synchronously by the corresponding
\texttt{End*} path. Its \texttt{NetworkSessionAction} now starts completed instead of waiting
for an event that cannot occur. A separate-process harness exercises the initialized dispatcher
under a ten-second watchdog because its process-lifetime statics have no test reset hook. The
harness also checks the four local identities, sign-in events, achievement completion, and
reinitialization cleanup. Chapter~\ref{ch:networking} covers the network side.
\end{historynote}

\section{Evidence and user implications}

The pinned headers and implementations establish the API shapes, fixed identities, file
layout, overlay hooks, and stubs. Unit tests cover persistence behavior and Guide state;
the isolated harness covers the dispatcher interaction. These tests do not establish
interoperation with Xbox Live, because CNA has no such transport.

\begin{center}
\begin{tabularx}{\textwidth}{@{}lX@{}}
\toprule
\textbf{Feature} & \textbf{Pinned-revision status} \\
\midrule
Local gamers & Four fixed identities; implemented in memory \\
Achievements & Locally persisted; catalog metadata and pictures absent \\
Leaderboards & Locally persisted with CNA-defined sorting and paging \\
Presence, privileges, profile & Local or synthetic state; no service enforcement \\
Friends and remote gamer lookup & Empty or unsupported \\
Most \texttt{Guide.Show*} calls & No-op stubs \\
Message-box and keyboard input & CNAEXT-rendered overlays \\
Avatar base API & Inert; see Chapter~\ref{ch:avatar} \\
\bottomrule
\end{tabularx}
\end{center}
